\section{Conclusions}
\label{sec:conclusion}

This study presented a morphology-aware validation framework for DRL-based battery dispatch in a distribution-network microgrid. The central result is conditional by design: the main tables evaluate six P4-selected candidates that pass the pre-declared morphology gates, rather than estimating unconditional SAC training reliability. Within this selected set, the simple-EBM objective cost is $11.788 \pm 0.120$\,M\textyen{}, and replay under Thevenin loss-only, Rint-only, and full PBM battery models changes the mean objective by less than 0.5\% relative to the simple model. This cross-fidelity result should be read as limited replay sensitivity under the tested IEEE 33-bus setting, not as a general robustness certificate. The more important observation is behavioral: economically similar checkpoints exhibit distinct SOC morphology, price-response patterns, raw-action infeasibility, and teacher-gap profiles. The reported P4 full-year evaluation records zero shield activation, while the normalized inventory-teacher action gap remains $0.301 \pm 0.087$ and infeasible-action dwell remains $6.1 \pm 7.3$\%, showing that shield execution, teacher alignment, and policy autonomy are separable diagnostics. The protocol therefore acts as an audit layer around SAC rather than as a new optimizer: it can select and rank candidates from a training pool, but it cannot create reliable policies from insufficient candidate diversity. This distinction motivates multi-seed reporting, consistent with DRL evaluation concerns raised in \cite{agarwal2021statistical}. The current evidence remains bounded by a single feeder topology, fixed BESS/PV placement, perfect-foresight oracle normalization, and a network audit that reports power-flow failures but not full voltage-margin, line-loading-margin, feeder-loss, or sensitivity distributions. Operation under forecast error, initial-SOC perturbation, abnormal weather, rolling-horizon execution, and changing BESS availability remains future work, consistent with EPSR studies that treat BESS and cloud-energy-storage operation under data-driven uncertainty as separate problems\cite{parvar2022epsrDRO,saini2024epsrUncertainty}. Overall, morphology-aware validation provides a compact engineering screen for deciding whether a learned battery policy merits further evaluation, additional training, or replacement by another candidate.
